\section{Carte PC--Fock complète, mémoire et stockage terminal}
\label{sec:ns-pcf-prospective-construction}

Nous fixons \(s>5/2\), \(\nu>0\), une coupure de Galerkin \(M\), une profondeur nonadique \(J\) et un horizon \(T<\infty\). Soit \(V_M\subset H^s_\sigma(\mathbb T^3)\) l'espace solénoïdal de moyenne nulle, et soient
\begin{equation}
 A_M=P_M(-\Delta)P_M,
 \qquad
 N_M(u)=P_M\mathbb P(u\cdot\nabla u).
 \label{eq:ns-import-galerkin-operators}
\end{equation}
Cette section construit le support PC--Fock complet, ses lecteurs, mémoires, retenue, stockage terminal et retour, qui interviennent dans la réalisation de référence de la section suivante. Son KYP local, son observabilité NS--OBS et les dominations LRDN--LCN sont exclusivement un \texttt{CONTRÔLE\_NON\_DÉTERMINANT} : ils ne conditionnent ni la métrique \(U^{\rm own}\), ni sa télescopie, ni son budget racine, ni la conclusion globale.

\subsection{Tour complète et lecteurs de degrés un et deux}

Le relèvement physique est pris dans l'espace de Fock symétrique complet :
\begin{equation}
 \mathcal V_M^{\rm PCF}(u)
 :=\bigoplus_{n\ge0}\frac{R_{\rm F}^n}{\sqrt{n!}}u^{\odot n},
 \qquad
 \|\mathcal V_M^{\rm PCF}(u)\|^2
 =e^{R_{\rm F}^2\|u\|_{H^s}^2}.
 \label{eq:ns-import-pcf-lift}
\end{equation}
La tour n'est pas tronquée. Pour \(z_{n,M}=u_M^{\odot n}\), la règle de Leibniz donne la récurrence exacte
\begin{equation}
 \dot z_{n,M}
 =L_{n,M}z_{n,M}
 +\mathcal N_{n,n+1;M}z_{n+1,M}
 +\mathcal F_{n,n-1;M}(f_M)z_{n-1,M},
 \qquad n\ge1,
 \label{eq:ns-import-carleman-all-degrees}
\end{equation}
où, sur la diagonale cohérente,
\begin{align}
 L_{n,M}u^{\odot n}
 &=-\nu n(A_Mu)\odot u^{\odot(n-1)}, \notag\\
 \mathcal N_{n,n+1;M}u^{\odot(n+1)}
 &=-nN_M(u)\odot u^{\odot(n-1)}, \notag\\
 \mathcal F_{n,n-1;M}(f)u^{\odot(n-1)}
 &=nf\odot u^{\odot(n-1)}.
 \label{eq:ns-import-carleman-blocks}
\end{align}

L'espace d'état est la somme orthogonale
\begin{equation}
 \mathscr K_{M,j}^{\rm PCF}
 =\mathscr K_{M,j}^{\rm HMT}\widehat\otimes\Gamma_s(V_M)
 \oplus\mathscr K_{M,j}^{\kappa}
 \oplus\mathscr K_{M,j}^{\rm micro}
 \oplus\mathscr K_{M,j}^{\rm sh},
 \label{eq:ns-import-full-carrier}
\end{equation}
où \(\mathscr K_{M,j}^{\rm HMT}\) conserve feuillet, orientation, \(q_A,q_V,T\), bord, parcours et les deux mémoires \(D/H\). Les degrés un et deux sont lus au moyen de
\begin{equation}
 \begin{aligned}
 z_M(u)&=\nu^{1/2}A_M^{1/2}\Lambda^su,\\
 w_M(u)&=\nu^{-1/2}A_M^{-1/2}\Lambda^sN_M(u),\\
 \mathcal Y_{\beta,M}\mathcal V_M^{\rm PCF}(u)
 &=\sqrt\beta\,z_M(u)
   \oplus\frac{1}{2\sqrt\beta}\,w_M(u).
 \end{aligned}
 \label{eq:ns-import-degree-one-two-reader}
\end{equation}
En conséquence,
\begin{equation}
 |B_{s,M}(u)|
 \le
 \|\mathcal Y_{\beta,M}\mathcal V_M^{\rm PCF}(u)\|^2
 =\beta\nu D_{s,M}(u)+\frac{\|w_M(u)\|^2}{4\beta}.
 \label{eq:ns-import-young-port}
\end{equation}
La seconde coordonnée est un opérateur linéaire uniformément borné sur \(\operatorname{Sym}^2H^s_\sigma\) :
\begin{equation}
 \sup_M\|\mathfrak C_M\|
 \le\widetilde M_{s,\beta,\nu},
 \qquad
 C_M^{\rm obs}:=\sqrt2\,\mathfrak C_M\pi_2,
 \qquad
 \Sigma_M:=\bigl(I+(C_M^{\rm obs})^*C_M^{\rm obs}\bigr)^{-1}
 \succeq\frac{I}{1+2\widetilde M_{s,\beta,\nu}^2}.
 \label{eq:ns-import-uniform-graph}
\end{equation}

\subsection{Mémoire bidirectionnelle et stockage terminal}

Avec
\begin{equation}
 r=\frac{\pi_{\rm HMT}}{\varphi_{\rm HMT}^2},
 \qquad
 D_{M,j}^{\rm mem}=(r-1)(M_{M,j}^+-M_{M,j}^-),
 \qquad
 H_{M,j}=(1+r)(M_{M,j}^++M_{M,j}^-),
 \label{eq:ns-import-dh}
\end{equation}
les deux mémoires satisfont
\begin{equation}
 \dot D_{M,j}^{\rm mem}=B_{s,M,j},
 \qquad
 \dot H_{M,j}=\frac{1+r}{r-1}|B_{s,M,j}|.
 \label{eq:ns-import-dh-dot}
\end{equation}
En posant
\[
 \delta_r=\frac{r-1}{r+1},
 \qquad
 \mathcal N_{M,j}:=\delta_rH_{M,j}-D_{M,j}^{\rm mem}
 =2(r-1)M_{M,j}^-\ge0,
\]
on obtient
\begin{equation}
 \dot{\mathcal N}_{M,j}=|B_{s,M,j}|-B_{s,M,j}.
 \label{eq:ns-import-negative-memory}
\end{equation}
Le stockage terminal explicite est
\begin{equation}
 \begin{aligned}
 G_{M,j,T}(t)
 &=E_{s,M,j}(t)
   +2(r-1)\bigl(M_{M,j}^-(T)-M_{M,j}^-(t)\bigr),\\
 G_{M,j,T}(t)&\ge E_{s,M,j}(t),\\
 \dot G_{M,j,T}+\nu D_{s,M,j}+|B_{s,M,j}|
 &=F_{s,M,j}.
 \end{aligned}
 \label{eq:ns-import-terminal-storage}
\end{equation}
Si
\[
 h^-_{M,j;t,T}(\tau)
 :=\mathbf1_{[t,T]}(\tau)\sqrt{\dot M^-_{M,j}(\tau)},
\]
la colonne
\begin{equation}
 \mathcal E_{M,j,T}^{\rm term}(t)\Xi
 :=2^{-1/2}\Lambda^su_{M,j}(t)
 \oplus\sqrt{2(r-1)}\,h^-_{M,j;t,T}
 \label{eq:ns-import-terminal-column}
\end{equation}
réalise le stockage comme un Gram :
\begin{equation}
 \|\mathcal E_{M,j,T}^{\rm term}(t)\Xi\|^2
 =G_{M,j,T}(t).
 \label{eq:ns-import-terminal-gram}
\end{equation}

La retenue commune entre profondeurs est
\begin{equation}
 \kappa_{M,j}
 =\frac12\left(E_j|B_{s,M,j+1}|-|B_{s,M,j}|\right)\ge0,
 \qquad
 E_jB_{s,M,j+1}=B_{s,M,j}.
 \label{eq:ns-import-carry}
\end{equation}
Sa coordonnée de sortie est fixée par
\begin{equation}
 q_{M,j}^{\kappa}
 :=\left(\frac{2(1+r)}{r-1}\kappa_{M,j}\right)^{1/2},
 \qquad
 \|q_{M,j}^{\kappa}\|^2
 =\frac{2(1+r)}{r-1}\kappa_{M,j}.
 \label{eq:ns-import-carry-port}
\end{equation}
Parallèlement,
\begin{equation}
 E_j\Delta_T\mathcal N_{M,j+1}
 =\Delta_T\mathcal N_{M,j}+2K_{M,j}(T),
 \qquad
 K_{M,j}(T):=\int_0^T\kappa_{M,j}(t)\,dt.
 \label{eq:ns-import-carry-telescope}
\end{equation}
Les innovations microscopiques et les bandes de Galerkin restent orthogonales :
\begin{equation}
 Q_{\rm micro}^{\rm APP}=I-J_9E_9,
 \qquad E_9y^{\rm micro}=0,
 \qquad \mathsf K_{\rm carry}y^{\rm micro}=0,
 \label{eq:ns-import-micro}
\end{equation}
et, degré par degré,
\begin{equation}
 \mathcal N_{n,n+1;3M}J_M^{(n+1)}
 =J_M^{(n)}\mathcal N_{n,n+1;M}
 +\mathcal S_{n,M}^{\rm sh},
 \qquad
 E_M^{(n)}\mathcal S_{n,M}^{\rm sh}=0.
 \label{eq:ns-import-shell}
\end{equation}

\subsection{Contrôle KYP fini non déterminant et cascade de la carte}

La force et son terme direct sont incorporés au moyen de
\begin{equation}
 \begin{aligned}
 x_{M,j+1}
 &=\mathcal A_{M,j}^{[j]}x_{M,j}
   +\mathcal B_{M,j,T}^{F,[j]}g_{M,j},\\
 y_{M,j}^{\rm loss}
 &=\mathcal C_{M,j,T}^{[j]}x_{M,j}
   +\mathcal D_{M,j,T}^{F,[j]}g_{M,j},
 \end{aligned}
 \label{eq:ns-import-affine-step}
\end{equation}
avec la décomposition orthogonale
\begin{equation}
 \mathcal C_{M,j,T}
 =\mathcal C_{M,j,T}^{(1)}
 \oplus\mathcal C_{M,j,T}^{(2)}
 \oplus\mathcal C_{M,j,T}^{\kappa}
 \oplus\mathcal C_{M,j,T}^{\rm micro}
 \oplus\mathcal C_{M,j,T}^{\rm sh}.
 \label{eq:ns-import-loss-decomposition}
\end{equation}
Le terme direct apparaît déjà dans
\[
 g_M=A_M^{-1/2}\Lambda^sP_Mf,
 \qquad
 a_M=\frac{g_M}{2\sqrt\nu},
 \qquad
 y_{D,M}=\sqrt\nu A_M^{1/2}\Lambda^su_M-a_M,
\]
pour lequel
\begin{equation}
 dG_{M,j,T}
 +\bigl(\|y_{D,M}\|^2+|B_{s,M,j}|\bigr)\,dt
 =\|a_M\|^2\,dt.
 \label{eq:ns-import-force-feedthrough}
\end{equation}

L'affinité est homogénéisée sans dupliquer la donnée. Soit
\[
 \mathscr G_M^{F,[j]}
 :=\bigoplus_{k=j}^{J-1}L^2(\Omega_k;\mathscr G^F_{M,k,T}),
 \qquad
 \mathbf x_{M,j}:=x_{M,j}\oplus(g_{M,j},\ldots,g_{M,J-1}),
\]
et soient \(e_j\) l'évaluation de la première entrée et \(\mathsf S_j\) le décalage de la queue. Nous définissons
\begin{equation}
 \mathbf\Gamma_{M,j}
 :=\begin{pmatrix}
  \mathcal A_{M,j}^{[j]}&\mathcal B_{M,j,T}^{F,[j]}e_j\\
  0&\mathsf S_j
 \end{pmatrix},
 \qquad
 \mathbf L_{M,j,T}
 :=\mathcal Q_{M,j,T}^{\rm loss,inc}
 \begin{pmatrix}
  \mathcal C_{M,j,T}^{[j]}&\mathcal D_{M,j,T}^{F,[j]}e_j
 \end{pmatrix}.
 \label{eq:ns-import-homogeneous-affine-step}
\end{equation}
Ici, \(\mathcal Q_{M,0,T}^{\rm loss,inc}=I\) ; pour \(j>0\), il supprime uniquement la copie héritée des canaux PC--Fock et agit comme l'identité sur la retenue, l'innovation microscopique et la bande spectrale.

\begin{theorem}[Contrôle de factorisation KYP finie non circulaire]
\label{thm:ns-import-affine-stein}
Au niveau terminal, on définit
\begin{equation}
 \mathbf U_{M,J,T}(t)
 :=(\mathcal E_{M,J,T}^{\rm term}(t))^*
   \mathcal E_{M,J,T}^{\rm term}(t).
 \label{eq:ns-import-terminal-metric}
\end{equation}
Une fois \(M,J,T\) fixés, soient \(A,B,C,D\) les quatre blocs du pas affine et \(U_+=\widetilde{\mathbf U}_{M,j+1,T}(t)\). Définissons
\begin{equation}
 Q:=I-B^*U_+B-D^*D,
 \qquad L:=B^*U_+A+D^*C.
 \label{eq:ns-import-kyp-q-l}
\end{equation}
Si \(Q\succ0\), la récurrence non circulaire est
\begin{equation}
 U:=A^*U_+A+C^*C+L^*Q^{-1}L.
 \label{eq:ns-import-affine-stein}
\end{equation}
Alors la matrice KYP complète satisfait
\begin{equation}
 \begin{aligned}
 \mathfrak S_{M,j,T}
 &:=
 \begin{pmatrix}
 U-A^*U_+A-C^*C&-A^*U_+B-C^*D\\
 -B^*U_+A-D^*C&I-B^*U_+B-D^*D
 \end{pmatrix}\\
 &=
 \begin{pmatrix}Q^{-1/2}L&-Q^{1/2}\end{pmatrix}^{*}
 \begin{pmatrix}Q^{-1/2}L&-Q^{1/2}\end{pmatrix}
 \succeq0.
 \end{aligned}
 \label{eq:ns-import-affine-supply}
\end{equation}
Pour \(Q\succeq0\), la variante avec pseudo-inverse est valable si \(\operatorname{Ran}L\subseteq\operatorname{Ran}Q^{1/2}\). Cette construction est finie : elle n'affirme pas que \(Q^{-1}\), \(U\) ou le gain restent uniformes lorsque \(M,J\) est retiré.
\end{theorem}

\begin{proof}
Par \eqref{eq:ns-import-affine-stein}, le bloc supérieur gauche est \(L^*Q^{-1}L\) ; les blocs non diagonaux sont \(-L^*\) et \(-L\), et le bloc inférieur droit est \(Q\). La multiplication de la ligne exposée dans \eqref{eq:ns-import-affine-supply} donne exactement les quatre blocs. La positivité précède ainsi l'extraction de toute racine.
\end{proof}

Ce théorème vérifie uniquement la carte finie lorsque ses propres hypothèses sont satisfaites. Il ne fournit aucune prémisse au développement de référence Stein--Lyapunov et ne limite pas l'uniformité démontrée par son Gram terminal et sa racine commune.

Dans les coordonnées homogénéisées, on prend \(\mathbf U_{M,j,T}=U\) et l'on note \(\mathbf R_{M,j,T}\) le relèvement de la ligne \(\begin{psmallmatrix}Q^{-1/2}L&-Q^{1/2}\end{psmallmatrix}\).
Avec cette convention, on retrouve, pour chaque troncature finie satisfaisant l'hypothèse KYP,
\[
 \mathbf U_{M,j,T}
 =\mathbf\Gamma_{M,j}^*\widetilde{\mathbf U}_{M,j+1,T}
  \mathbf\Gamma_{M,j}
 +\mathbf L_{M,j,T}^*\mathbf L_{M,j,T}
 +\mathbf R_{M,j,T}^*\mathbf R_{M,j,T}.
\]

La colonne
\begin{equation}
 [\mathbf x]_{\mathbf U_{M,j,T}}
 \longmapsto
 [\mathbf\Gamma_{M,j}\mathbf x]_{
  \widetilde{\mathbf U}_{M,j+1,T}}
 \oplus\mathbf L_{M,j,T}\mathbf x
 \oplus\mathbf R_{M,j,T}\mathbf x
 \label{eq:ns-import-local-isometry}
\end{equation}
est isométrique. Sa dilatation de Julia et sa composition produisent \(\mathfrak U_{M,J,T}\) avec
\begin{equation}
 \mathfrak U_{M,J,T}^{\sharp}\mathfrak U_{M,J,T}=I
 \label{eq:ns-import-cascade-isometry}
\end{equation}
et le télescope
\begin{equation}
 \begin{aligned}
 \|\mathbf x_{M,0}\|_{\mathbf U_{M,0,T}}^2
 ={}&\mathbb E_{0:J}
 \|\mathbf x_{M,J}\|_{\mathbf U_{M,J,T}}^2\\
 &+\sum_{j=0}^{J-1}\mathbb E_{0:j}
 \left(\|\mathbf L_{M,j,T}\mathbf x_{M,j}\|^2
       +\|\mathbf R_{M,j,T}\mathbf x_{M,j}\|^2\right).
 \end{aligned}
 \label{eq:ns-import-stein-telescope}
\end{equation}

\subsection{Budget de la carte et retour de la vitesse}

Sur la diagonale physique, la première coordonnée terminale est \(G_{M,J,T}\), et les autres coordonnées forment la somme des pertes :
\begin{equation}
 \|Z_{M,J,t}\|^2
 =\mathbb E_{0:J}G_{M,J,T}(t)
 +\sum_{j=0}^{J-1}\mathbb E_{0:j}\Lambda_{M,j}([0,t]).
 \label{eq:ns-import-ledger}
\end{equation}
Puisque \(q_{M,j}^{\kappa}\) est l'une de ces coordonnées, \eqref{eq:ns-import-carry-telescope} n'est intégré qu'une seule fois :
\begin{equation}
 2\sum_{j=0}^{J-1}\mathbb E_{0:j}K_{M,j}(T)
 \le\|\Xi_{M,0}\|^2+Q_f(T).
 \label{eq:ns-import-carry-funded}
\end{equation}
D'où,
\begin{equation}
 \mathbb E_{0:J}G_{M,J,T}(0)
 \le E_{s,M}(P_Mu_0)
 +\Delta_T\mathcal N_0+\|\Xi_{M,0}\|^2+Q_f(T).
 \label{eq:ns-import-root-bound}
\end{equation}
La racine est commune à toutes les coupures :
\begin{equation}
 \Xi_{M,0}=J_{0\to M}\Xi_0,
 \qquad J_{0\to M}^*J_{0\to M}=I,
 \qquad \mathcal N_{M,0}=\mathcal N_0.
 \label{eq:ns-import-common-root}
\end{equation}
De plus,
\begin{equation}
 \|\Xi_{M,0}\|^2
 \le C_{\rm HMT}
 +(1+2\widetilde M_{s,\beta,\nu}^2)
 e^{R_{\rm F}^2\|u_0\|_{H^s}^2},
 \qquad
 Q_f(T)\le C_{\nu,s}\int_0^T\|f(t)\|_{H^{s-1}}^2\,dt.
 \label{eq:ns-import-root-data-bound}
\end{equation}
Par conséquent, pour cette carte et ces \(M,J,T\), nous posons
\begin{align}
 \mathcal B_{M,J,T}^{\rm route}:={}&
 \frac12\|u_0\|_{H^s}^2
 +|\Delta_T\mathcal N_{M,0}|+C_{\rm HMT}
 +(1+2\widetilde M_{s,\beta,\nu}^2)
 e^{R_{\rm F}^2\|u_0\|_{H^s}^2}\notag\\
 &+C_{\nu,s}\int_0^T\|f(t)\|_{H^{s-1}}^2\,dt
 <\infty
 \label{eq:ns-import-explicit-budget}
\end{align}
La présence explicite de \(\Delta_T\mathcal N_{M,0}\) empêche d'utiliser cette expression comme preuve indépendante d'uniformité. C'est un diagnostic local de la carte ; il ne remplace pas le budget de référence \eqref{eq:nsclose-owner-data-budget}. L'identité terminale et l'inégalité de Young donnent, pour la coupure fixée,
\begin{equation}
 \boxed{
 \sup_{0\le t\le T}\|Z_{M,J,t}\|^2
 +\frac{\nu}{2}\int_0^TD_{s,M}(t)\,dt
 \le\mathcal B_{M,J,T}^{\rm route}.}
 \label{eq:ns-import-uniform-budget}
\end{equation}

L'expression de la vitesse est normalisée par
\begin{equation}
 \iota_{s,M}u=2^{-1/2}\Lambda^su,
 \qquad
 \mathcal R_{1,M}=\sqrt2\,\Lambda^{-s}\pi_1,
 \qquad
 \mathcal R_{1,M}\iota_{s,M}=I_{V_M}.
 \label{eq:ns-import-grade-one-return}
\end{equation}
Le lecteur de la cascade est
\begin{equation}
 \mathcal R_{M,J,T}^{\rm phys}
 :=\mathcal R_{1,M}\mathfrak U_{M,J,T}^{\sharp}.
 \label{eq:ns-import-full-return}
\end{equation}
Par \eqref{eq:ns-import-cascade-isometry},
\begin{equation}
 \boxed{
 \mathcal R_{M,J,T}^{\rm phys}
 \mathfrak U_{M,J,T}\mathcal V_M^{\rm PCF}(u_M(t))
 =u_M(t).}
 \label{eq:ns-import-exact-return}
\end{equation}
Le retour utilise l'adjoint métrique de la même cascade et retrouve l'expression hydrodynamique sans sélectionner de branche probabiliste.
